\documentclass[a4paper,8pt,twoside]{article}
\usepackage[utf8]{inputenc}
\usepackage[T1]{fontenc}
\usepackage[french]{babel}
\usepackage{import}

\import{C:/Users/lasca/Desktop/Cours de mathématiques/(Modèle de cours et d'exercices)/Autres types de documents (khôlles et autres)/Plans d'agreg}{metaplans.tex}
\import{C:/Users/lasca/Desktop/Cours de mathématiques/(Modèle de cours et d'exercices)}{preambulestandard.tex}
\import{C:/Users/lasca/Desktop/Cours de mathématiques/(Modèle de cours et d'exercices)}{environnementspersonnels.tex}
\import{C:/Users/lasca/Desktop/Cours de mathématiques/(Modèle de cours et d'exercices)}{commandespersonnelles.tex}

\newcommand{\showcorrection}{Commenter cette ligne pour faire disparaître les corrections des exercices}

\sectionfont{\fontsize{13}{15}\selectfont}
\setlength\parindent{0pt}
\linespread{0.9}

\definecolor{amber}{rgb}{1.0, 0.49, 0.0}
\hypersetup{
    linkcolor=amber,
}

\newcommand{\titreducoursmod}[1]{
\titreducours{#1}
\addtocontents
  {toc}
  {\protect\contentsline
             {section}
             {\protect\numberline{}#1}{}{}}
\setcounter{section}{0}
}

\begin{document}

%Commandes de dernière minute
\makeatletter
\let\ref\@refstar
\makeatother

\begin{landscape}

\thispagestyle{empty}

\title{\huge Un peu plus de 69 métaplans pittoresques... \\ pour se tenir à jour}
\date{\today}
\author{Louis Lascaud}

\thispagestyle{empty}

\maketitle

\thispagestyle{empty}

\begin{multicols*}{2}

\thispagestyle{empty}


\vspace{1.4 cm}

\tableofcontents
\addtocontents{toc}{\protect\thispagestyle{empty}}
\vspace{1 cm}

Cette feuille complète le document « 40 développements gouleyants » en rajoutant des développements et en actualisant le couplage pour qu'il corresponde à l'évolution du programme pour la session en cours. Le nouveau couplage apparaît dans le tableau suivant : \url{https://louis-lascaud.github.io/l-lscd-mathematiques/Docs/depotprive/New%20couplage.xlsx} à télécharger au format Excel. \\ \vspace{0.3 cm}

\end{multicols*}
\end{landscape}

\newpage

\titreducoursmod{(202) Nombres complexes de partie réelle un demi. Droites critiques. Exemples et applications}

\section{Section 1}

\end{multicols*}
\end{landscape}

\end{document}